\documentclass[11pt]{article}
\usepackage[margin=1in]{geometry}
\usepackage{amsmath,amssymb,amsthm,bm}
\usepackage{mathtools}
\usepackage{enumitem}

\title{Theory Note: Gaussian-Process ROM Used in Inverse Workflows}
\author{romtools/workflows/inverse}
\date{\today}

\begin{document}
\maketitle

\section{Purpose and Scope}
This note documents the mathematical model implemented by the lightweight Gaussian-process (GP)
surrogate in \texttt{romtools/rom/qoi\_surrogates.py}, as used by
\texttt{mf\_eki\_drivers.py} and \texttt{mf\_vi\_drivers.py}.
The focus is the exact training and prediction equations represented in code, including:
\begin{itemize}[leftmargin=1.5em]
\item scalar-output GP regression,
\item vector-output reduction via POD and one GP per POD coefficient,
\item optional parameter and target normalization,
\item optional grid-search kernel tuning by log marginal likelihood.
\end{itemize}

\section{Data and Notation}
Let the training inputs be
\[
X = \begin{bmatrix}
\bm{x}_1^T\\
\vdots\\
\bm{x}_n^T
\end{bmatrix} \in \mathbb{R}^{n\times d},
\]
where $n$ is the number of training samples and $d$ is parameter dimension.

For scalar QoI, the targets are $y\in\mathbb{R}^n$.
For vector QoI with physical dimension $m$, the targets are
\[
Y=\begin{bmatrix}
\bm{y}_1^T\\
\vdots\\
\bm{y}_n^T
\end{bmatrix}\in\mathbb{R}^{n\times m}.
\]

\section{Kernel and GP Prior}
The implementation uses an isotropic squared-exponential kernel:
\[
k(\bm{x},\bm{x}') = \sigma_f^2
\exp\!\left(-\frac{\|\bm{x}-\bm{x}'\|_2^2}{2\ell^2}\right),
\]
with length scale $\ell>0$ and signal variance $\sigma_f^2>0$.

Define
\[
K(X,X)_{ij}=k(\bm{x}_i,\bm{x}_j).
\]
Training adds independent Gaussian noise plus fixed jitter:
\[
\widetilde{K}=K(X,X)+(\sigma_n^2+\varepsilon)I_n,
\]
where $\sigma_n^2$ is the configured noise variance and $\varepsilon$ is jitter.

\section{Scalar-Output GP Training}
For scalar targets $y\in\mathbb{R}^n$, the code computes a Cholesky factor
\[
\widetilde{K}=LL^T,
\]
then solves
\[
\bm{\alpha}
=\widetilde{K}^{-1}y
=L^{-T}L^{-1}y.
\]
The stored model state is $(X,\bm{\alpha},L)$.

\subsection*{Prediction}
For query points $X_\star\in\mathbb{R}^{n_\star\times d}$:
\[
K_\star = K(X_\star,X)\in\mathbb{R}^{n_\star\times n},\qquad
\hat{y}_\star = K_\star \bm{\alpha}.
\]
The implementation returns the posterior mean only (no posterior variance path is exposed).

\section{Log Marginal Likelihood for Kernel Selection}
If hyperparameter tuning is enabled, candidate kernels are evaluated using:
\[
\log p(y\mid X,\theta,\sigma_n^2)
=
-\frac{1}{2}y^T\widetilde{K}^{-1}y
-\frac{1}{2}\log\det(\widetilde{K})
-\frac{n}{2}\log(2\pi),
\]
with $\theta=(\ell,\sigma_f^2)$.

Using Cholesky $\widetilde{K}=LL^T$:
\[
\log\det(\widetilde{K}) = 2\sum_{i=1}^n \log L_{ii},
\]
so the implemented expression is
\[
-\frac{1}{2}y^T\bm{\alpha} - \sum_{i=1}^n\log L_{ii} -\frac{n}{2}\log(2\pi).
\]
Candidates that fail Cholesky factorization are assigned $-\infty$.
The best candidate on the configured grid is selected.

\section{Vector QoI: POD + Per-Coefficient GPs}
For $Y\in\mathbb{R}^{n\times m}$:
\begin{enumerate}[leftmargin=1.5em]
\item Compute mean
\[
\bar{\bm{y}}=\frac{1}{n}\sum_{i=1}^n \bm{y}_i.
\]
\item Center data: $Y_c=Y-\mathbf{1}\bar{\bm{y}}^T$.
\item Compute thin SVD:
\[
Y_c = U\Sigma V^T.
\]
\item Choose POD rank $k$ by cumulative singular-value energy:
\[
\frac{\sum_{j=1}^k \sigma_j^2}{\sum_{j}\sigma_j^2}
\geq \eta,
\]
where $\eta=$ \texttt{pod\_energy\_fraction}, then apply optional cap
\texttt{max\_pod\_modes}; enforce $k\ge 1$.
\item Define POD modes $\Phi = V_{1:k}\in\mathbb{R}^{m\times k}$ and coefficients
\[
C = Y_c\Phi \in \mathbb{R}^{n\times k}.
\]
\item Train one scalar GP per column $C_{:,j}$.
\end{enumerate}

For a query $\bm{x}_\star$, let predicted coefficients be $\hat{\bm{c}}_\star\in\mathbb{R}^k$.
Reconstruction is
\[
\hat{\bm{y}}_\star = \bar{\bm{y}} + \Phi\hat{\bm{c}}_\star.
\]

\section{Normalization Options}
\subsection*{Input normalization (\texttt{normalize\_parameters})}
For each input dimension:
\[
x'_{i,r}=\frac{x_{i,r}-x^{\min}_r}{x^{\max}_r-x^{\min}_r},
\]
with zero range replaced by $1$ to avoid division by zero.

\subsection*{Target normalization (\texttt{normalize\_targets})}
If enabled, each target channel (scalar output or each POD coefficient) uses min-max scaling:
\[
t'=\frac{t-t_{\min}}{t_{\max}-t_{\min}},
\]
again replacing zero scale by $1$.
Predictions are mapped back with the inverse affine transform.

\section{Noise-Variance Policy}
For each trained scalar GP target:
\begin{itemize}[leftmargin=1.5em]
\item If \texttt{noise\_variance} is provided, use it directly.
\item Else if \texttt{auto\_noise\_variance=True}, set
\[
\sigma_n^2 = \gamma \,\mathrm{Var}(t),
\]
where $\gamma=$ \texttt{noise\_variance\_fraction}.
\item Else use default $\sigma_n^2=10^{-10}$.
\end{itemize}
Jitter is a fixed additive term ($10^{-12}$ in the lite regressor).

\section{Computational Cost}
For one scalar GP with $n$ training samples:
\begin{itemize}[leftmargin=1.5em]
\item Kernel assembly: $\mathcal{O}(n^2 d)$.
\item Cholesky + solves: $\mathcal{O}(n^3)$.
\item Mean prediction at one query: $\mathcal{O}(n d)$ for kernel row + $\mathcal{O}(n)$ dot product.
\end{itemize}

For vector QoI with $k$ retained POD modes, training cost is approximately
$k$ times the scalar-GP training cost (plus one SVD of $Y_c$).
This is the primary reason training time can grow rapidly as training set size increases.

\section{Algorithm Summary}
\begin{enumerate}[leftmargin=1.5em]
\item Optionally normalize parameters.
\item If scalar QoI, train one GP on scalar targets.
\item If vector QoI:
\begin{enumerate}[label*=\arabic*.,leftmargin=1.5em]
\item compute POD basis from centered training QoIs,
\item project to POD coefficients,
\item optionally normalize coefficients,
\item train one GP per coefficient.
\end{enumerate}
\item At prediction time, evaluate all needed GP means and (for vector QoI) reconstruct through POD.
\end{enumerate}

\section{Implementation-Constrained Remarks}
The current implementation is intentionally lightweight:
\begin{itemize}[leftmargin=1.5em]
\item posterior mean is returned, but predictive variance is not exposed,
\item no gradient-based hyperparameter optimization is used; tuning is grid search only,
\item training uses direct dense linear algebra (Cholesky) and therefore scales cubically in sample count.
\end{itemize}

\end{document}
